\scanpage{30}
\mspage{26}
% archive sheet number: 29

On fera attention que cette notion de somme n'est \uline{pas} une notion de somme au sens catégorique (pour la notion de (co)morphisme de magasins introduite jusqu'à présent. \unsure{Ce ne serait pas normal|Ce qui serait peu normal}, vu l'interprétation géométrique et le renversement des flèches -- \unsure{intuitivement|intentionnel|interstices}, la somme de magasins correspond bien à une \uline{somme} d'espaces (et non à un produit, comme pourrait le suggérer le principe de renversement des flèches).
% slanted
\margin{vu la condition de compatibilité avec $|\circ|$}
% slanted
\margin{\deleted{\unsure{semble}}}

Si on part d'un magasin $\unsure{M|\mathcal M}$, on peut songer à définir ses composantes connexes $M_i$, comme les composantes connexes pour la relation $\ll$ : On trouve une décomposition en somme, à condition de supposer l'axiome suivant (qui sera \unsure{toujours|beaucoup} \unsure{confirmé|employé} par la suite)

\margin{\rule{0.35pt}{1.5cm}}
\uline{Mag conn$^{\unsure{0|\circ}}$}\quad Si $X,Y\in M$ appartiennent à des comp. connexes distinctes de $M$, on a
\[
X\mathrel{|\circ|}Y.
\]
